\chapter{Morphophonology}\label{ch:morphophonology}

\section{Morphological domains and grade selection}\label{sec:morphophon}

Morphological constructions select vocalic and consonantal grades at the
boundary where their exponents attach. Front and mixed grades select the
coronal and labial PAL correspondences of §\,\ref{sec:soft-coronals};
their separate velar specifications select the alternants described in
§\,\ref{sec:velar-fronting}. Antipassive sibilant replacement is a separate
operation. High front vowels do not impose a further categorical
sibilant mutation after morphology has selected the consonantal grade.

The nominal pair \ird{perek}/\ird{pérk-} combines loss of an
unstressed vowel with selection of the strong grade, whereas
\ird{lobek}/\ird{lúbk-} selects the raised strong grade. In the
principal spoken variety, strong and raised strong grades of short
mid-vowel roots both have long high-vowel outputs, while conservative
pronunciation preserves their height difference (§\,\ref{sec:vowel-raising}). The irrealis \ird{zďíčev},
from \ird{zděk-}, has a raised strong root vowel without vowel loss.
These distributions distinguish three independent
properties: membership in the fleeting-vowel class, selection of a raised grade,
and selection of a quantity grade. Each belongs to the lexeme and its
morphological construction; none adds a vowel phoneme to the inventory.

Within each construction, the full or bound stem is selected before its
exponent is attached and the associated consonant mutation applies:

\begin{center}
\textsc{construction and stem class} $\rightarrow$ \textsc{stem grade}
$\rightarrow$ \textsc{suffixation and consonant mutation} $\rightarrow$
\textsc{syllabification and licensed repair} $\rightarrow$
\textsc{coda vocalisation} $\rightarrow$ \textsc{voicing}
$\rightarrow$ \textsc{phonetic and orthographic realisation}
\end{center}

Lexical preverbs select no independent consonantal mutation, vowel-height
change, or quantity grade. Their derivatives retain the base's thematic class
and undergo the construction-selected alternations at subsequent suffixal
boundaries. Productive examples of \ird{le-}, \ird{ra-}, \ird{vra-},
\ird{ša-}, \ird{na-}, \ird{ru-}, and \ird{po-} are confined here to
consonant-initial bases, since the present description establishes no general
resolution of hiatus at these prefix--root boundaries. Combinations with
vowel-initial roots require individual lexical specification; automatic
contraction, deletion, epenthesis, and glide insertion cannot be inferred from
the prefix alone. Word-edge glottal insertion applies to the initial
vowel of the complete lexical prosodic word, rather than automatically to
an internal root edge (§\,\ref{sec:glottal-onset}); it consequently
supplies no general repair for these prefix--root combinations.

Consonant mutation and syllabification precede coda /l/ and /v/
vocalisation. The resulting [w] has sonorant behaviour when final devoicing
and cluster voicing apply to the remaining obstruents. Prevocalic onset /v/
is subsequently realised as [ʋ] or, where secondary palatal articulation is
licensed, [ʋʲ] (§§\,\ref{sec:l-vocalisation} and~\ref{sec:v-realisation}).

Thus \ird{lobek} selects \ird{lúbk-} before transitive \ird{-u}, while
\ird{perek} selects \ird{pérk-}; a front-vocalic ending subsequently yields
\ird{lúbcí} or \ird{pérce}. The stem \ird{zděk-} selects its raised modal
grade before \ird{-ěv}, whose velar mutation gives \ird{zďíčev}. Outer
anchoring morphology has its own stem-selection conditions
(§\,\ref{sec:nm-def-allomorph}).

Antipassive replacement or augmentation may accompany an independently
selected quantity grade. Its output supplies the input for subsequent
licensed suffixal boundaries and the ordinary phonology. Analogical
relations are encoded in the stored stem classes; they do not operate on
already assigned spellings.

\section{Consonantal grades and construction classes}

\subsection{Coronal and labial PAL correspondences}\label{sec:soft-coronals}

\textsc{pal} identifies a set of coronal and labial correspondences selected
by morphology. Membership in this set concerns the distribution of the
alternation, whereas the resulting articulation depends on the consonant.
The coronal stops and sonorants select [c\,ɟ\,ɲ\,ʎ], the alveolar sibilants
select postalveolars, and labials receive secondary palatal articulation.
These outcomes neither constitute one uniform phonetic operation nor require
an additional sibilant series.

\begin{table}[ht]
	\footnotesize
	\caption{Coronal correspondences selected by PAL.}\label{tab:soft-coronals}
	\begin{tabularx}{0.9\textwidth}{L L L L}
		\toprule
		{\sc input} & {\sc output} & {\sc spelling} & {\sc illustration} \\
		\midrule
		/t/ & [c] & tě / ť & \ird{rat} $\rightarrow$ \ird{ratě} \\
		/d/ & [ɟ] & dě / ď & \ird{rad} $\rightarrow$ \ird{radě} \\
		/n/ & [ɲ] & ně / ň & \ird{ran} $\rightarrow$ \ird{raně} \\
		/l/ & [ʎ] & lě / ľ & \ird{ral} $\rightarrow$ \ird{ralě} \\
		/s/ & /ʃ/ & š & schematic /s/ + \textsc{pal} $\rightarrow$ /ʃ/ \\
		/z/ & /ʒ/ & ž & schematic /z/ + \textsc{pal} $\rightarrow$ /ʒ/ \\
		/ts/ & /tʃ/ & č & schematic /ts/ + \textsc{pal} $\rightarrow$ /tʃ/ \\
		\bottomrule
	\end{tabularx}
\end{table}

The paired coronal spellings give the form before /ɛ/ first and the
consonant letter used elsewhere second. For /p\,b\,m\,v/, PAL supplies a secondary palatal specification while
retaining labial primary articulation. The outputs are represented
morphophonologically as /pʲ\,bʲ\,mʲ\,vʲ/; prevocalic realisation gives
[pʲ\,bʲ\,mʲ\,ʋʲ], since /v/ subsequently receives its approximant
realisation (§\,\ref{sec:v-realisation}). This notation identifies the
selected consonantal properties without positing an independently established
lexical contrast for every member of a soft labial series.

The standard soft-stop and lateral outputs are [c\,ɟ\,ʎ], rather than
merely [tʲ\,dʲ\,lʲ]. Their occurrence in a productive correspondence does
not alone establish lexical phonemic contrasts: the restricted lexical
distributions of /c, ɟ/ and the evidential limits for /ʎ/ are described in
§\,\ref{sec:consonant-inventory}. The postalveolars, like /ɲ/, have lexical
as well as morphologically derived occurrences.

PAL leaves already soft coronals and postalveolars unchanged. Rhotic /r/
has no categorical PAL alternant; any secondary [rʲ] realisation is
noncontrastive phonetic variation whose conditioning is not specified here.
Nor does PAL replace /r/ with learned \ird{ř} or /ʒ/. Lexical /kʲ/ is
independent of these correspondences. Plain /k/ has a separate,
class-selected velar output, so PAL membership and velar selection jointly
identify a consonantal grade (§\,\ref{sec:suffix-classes}).

Morphologically conditioned phonology can accompany an overt exponent
without itself constituting the sole expression of a category
\cite[ch.~3]{inkelas2014}. In the Distributed Morphology analysis used here,
the correspondences are stated as readjustments licensed by the selecting
exponent at a local boundary. This formal choice does not turn the
correspondence set into a single articulatory process.
The target need not belong to a bare Root: the linker \ird{-e} may act on
the exposed final consonant of an inflected complex. The precise
single-boundary mapping is defined in §\,\ref{sec:boundary-mapping}.

\subsection{Construction-specific grade selection}\label{sec:suffix-classes}

A construction selects a consonantal grade as well as a full or bound stem
and any vocalic allomorph. The consonantal grade specifies whether the
coronal and labial PAL correspondences apply and which output participating
plain /k/ receives. Only the four combinations in
Table~\ref{tab:licensed-grades} are licensed; the two properties are not
freely recombined into further grades.

\begin{table}[ht]
\footnotesize
\caption{Licensed combinations of PAL and velar selection.}\label{tab:licensed-grades}
\begin{tabularx}{0.95\textwidth}{L C L L}
\toprule
{\sc grade} & {\sc coronal/labial PAL} & {\sc input /k/} & {\sc other inputs} \\
\midrule
Plain & no & unchanged /k/ & unchanged \\
Front & yes & /ts/ & PAL correspondence where defined; otherwise unchanged \\
Mixed & yes & /tʃ/ & PAL correspondence where defined; otherwise unchanged \\
Restricted velar-front & no & /ts/ & unchanged \\
\bottomrule
\end{tabularx}
\end{table}

The shortened names \term{plain}, \term{front}, \term{mixed}, and
\term{restricted} refer to these four specifications. PAL identifies the
shared coronal and labial correspondences of the front and mixed grades,
rather than serving as another name for the front grade as a whole.
Thus restricted velar-front is not a weaker pronunciation of PAL: it
selects /ts/ from /k/ while preserving a boundary-final plain coronal.
The label \term{mixed} identifies the combination of PAL with velar
/tʃ/; it posits no pronounced or underlying /j/ in a suffix.

Membership belongs to an exponent in its construction, not to its initial
vowel letter. Czech provides a comparative case in which similar surface
suffix vowels accompany different velar outcomes
\cite{caviranivandenwyngaerd2026}; this supports the analytical separation
without determining Iridian's particular assignments or history.

\begin{table}[ht]
\footnotesize
\caption{Construction-specific selection of consonantal grades.}\label{tab:morph-trigger}
\begin{tabularx}{\textwidth}{L L}
\toprule
{\sc grade} & {\sc selecting constructions} \\
\midrule
Plain & indirect evidential \ird{-al}/\ird{-l}; causative \ird{-r-}, \ird{-er-}; conditional \ird{-ej}; nonmutating case exponents, including Class~I plural \ird{-e} and Class~V \ird{-e}; nonmutating aspect and converb exponents \\
Front & restricted causative \ird{-ěr-}; front-grade nominal exponents, including \ird{-í}, \ird{-e}, \ird{-ě}, \ird{-ení} and direct plural \ird{-a} in Class~VII; simultaneous \ird{-ěc}; linker \ird{-e}; stative adverb \ird{-ě}; emphatic \ird{-ěsk-}; denominative \ird{-ěvat-} \\
Mixed & irrealis \ird{-ěv}, with separately specified vocalic and cluster allomorphs where licensed \\
Restricted velar-front & negative irrealis \ird{-ir}; relational \ird{-(e)zn-} \\
\bottomrule
\end{tabularx}
\end{table}

An exponent can be represented as a pair of segmental material and selected
consonantal grade. Conventional suffix labels such as \ird{-ě} and
\ird{-ěv} identify morphemes whose surface spellings depend on the selected
output, rather than invariant strings of letters:

\begin{center}
\footnotesize
\begin{tabular}{llll}
\toprule
{\sc exponent} & {\sc segments} & {\sc grade} & {\sc construction} \\
\midrule
\ird{-e} & /ɛ/ & Front & linker \\
\ird{-ě} & /ɛ/ & Front & relevant nominal endings \\
\ird{-ěc} & /ɛts/ & Front & simultaneous converb \\
\ird{-ěv} & /ɛv/ & Mixed & irrealis \\
\ird{-ir} & /ɪr/ & Restricted & negative irrealis \\
\ird{-ej} & /ɛj/ & Plain & conditional \\
\bottomrule
\end{tabular}
\end{center}

Homophonous exponents need not have the same specification. Circumstantial
converb \ird{-ení} is plain-grade, whereas front-grade nominal
\ird{-ení} selects PAL and velar /ts/. Class~I plural \ird{-e} in
\ird{býle} and Class~V indirect \ird{-e} in \ird{ráne} are plain;
the homophonous indirect \ird{-e} of Class~II selects the front grade.
Conversely, direct plural \ird{-a} in Class~VII \ird{obľa} selects the
front grade despite lacking a front vowel. These assignments accompany
the nine declension classes in §\,\ref{sec:nm-declension}.

Causative \ird{-ěr-} selects the front grade only in its listed
mutating class, as in \ird{vad-ěr-ime} $\rightarrow$
\ird{vaďérime}; plain causative \ird{-r-} and \ird{-er-} leave
the preceding consonant unchanged (§\,\ref{sec:aspect-caus}). An
intervening applicative or antipassive blocks selection of the root:
the causative targets only its immediately preceding morphological boundary.

Vocalic selection remains independent: \ird{perek} selects bound
\ird{pérk-}, \ird{lobek} selects \ird{lúbk-}, and irrealis
\ird{zděk-} selects a raised strong vowel separately from its mixed-grade
consonant. Linker \ird{-e} and participant nominaliser \ird{-ou}
select perfective bound \ird{-ik-}; only the linker additionally selects
the front consonantal grade. Antipassive stem replacement or augmentation
belongs to a distinct construction (§\,\ref{sec:sibilant-replacement}).

A new root may enter an established inflectional or derivational class,
but a suffix vowel alone assigns neither an exceptional vocalic grade
nor an unlisted cluster allomorph.

\subsection{One mapping at each selected boundary}\label{sec:boundary-mapping}

For a boundary licensed by a morphological construction, let \(g\) be its
selected grade and let \(C\) be the final consonant of the selected input.
The input may already contain derivational or inflectional morphology.
Here \(C\) belongs to the selected input, not to the exponent about to
attach; the mapped stem combines with that exponent's independently
specified overt material. Let \(\kappa=1\) identify an input whose final
/k/ belongs to the participating velar class; otherwise \(\kappa=0\).
Write \(P(C)\) for the coronal or labial PAL correspondence in
§\,\ref{sec:soft-coronals}, and \(K_g\) for the output assigned to plain
/k/ by \(g\). The consonantal mapping is:

\[
 M_{g,\kappa}(SC)=S\,f_{g,\kappa}(C).
\]
\[
 f_{g,\kappa}(C)=
 \begin{cases}
 K_g, & C=\text{/k/}\text{ and }\kappa=1,\\
 P(C), & \text{PAL is selected and }C\text{ belongs to its correspondence set},\\
 C, & \text{otherwise}.
 \end{cases}
\]

The alternatives inspect the original boundary-final consonant. A
/k/-input in the participating class therefore receives the grade's velar
output once; the resulting
/ts/ is not submitted to \(P\) at that same boundary. An original /ts/,
by contrast, receives /tʃ/ in either grade with PAL. The distinction is
visible without assigning lexical meanings to the following schematic
inputs:

\begin{table}[ht]
\footnotesize
\caption{Outputs from participating velars and coronal inputs under each licensed grade.}\label{tab:boundary-map}
\begin{tabularx}{0.85\textwidth}{L C C C C}
\toprule
{\sc input} & {\sc plain} & {\sc front} & {\sc mixed} & {\sc restricted} \\
\midrule
/k/ (\(\kappa=1\)) & /k/ & /ts/ & /tʃ/ & /ts/ \\
/ts/ & /ts/ & /tʃ/ & /tʃ/ & /ts/ \\
/t/ & /t/ & [c] & [c] & /t/ \\
\bottomrule
\end{tabularx}
\end{table}

A nonparticipating /k/ remains unchanged under each grade; choosing a
grade does not itself assign an input to the participating velar class.
A boundary licenses only this one application. It cannot be reused to
mutate its own output. A distinct outer exponent can license another
application, whose input is then the exposed edge of the complex selected
by that exponent. For example, free perfective \ird{pjášček} first
supplies bound \ird{pjáščik-}; the linker's front-grade boundary maps
its final /k/ to /ts/, giving \ird{pjáščice}. It neither changes the
earlier /tʃ/ nor returns to the original root.

The mapping does not search leftwards through a cluster or across an
intervening suffix. If the selected input ends in a vowel, this
consonant-targeting operation has no target; any linked vowel replacement,
contraction, or extra consonant belongs to independently specified
allomorphy. Listed \ird{pjašt-}/\ird{pjášč-} in the irrealis likewise
supplies a construction-specific cluster stem, rather than a productive
/t/ $\rightarrow$ /tʃ/ exception added to \(P\). Morphological selection
establishes such an input before the boundary mapping and ordinary
phonetic realisation apply.

Orthography is assigned to the resulting consonantal structure. The letter
\orth{c} for /ts/ must not be confused with IPA [c], written \orth{ť},
and the \orth{j} of a labial digraph supplies no independent segment
to this mapping.

\section{Vocalic alternations}

\subsection{Thematic nominal endings}\label{sec:thematic-vowels}

The Class~VIII forms \ird{bro}, \ird{bra}, and \ird{brum} contain a constant consonantal
root \ird{br-} and the class-specific endings \ird{-o}, \ird{-a}, and
\ird{-um}. Their terminal vowels realise the case paradigm rather than
root-vowel ablaut:

\pex
\a \ird{br-} + \ird{-o} $\rightarrow$ \ird{bro}
\a \ird{br-} + \ird{-a} $\rightarrow$ \ird{bra}
\a \ird{br-} + \ird{-um} $\rightarrow$ \ird{brum}
\xe

In the bare-root Class~V represented by \ird{rán}, the singular endings are
zero, \ird{-ou}, and \ird{-e}. The long root vowel remains constant in
\ird{rán}, \ird{ránou}, and \ird{ráne}. The ending alternation leaves the root vowel and its quantity unchanged.

The anchored direct forms of Class~VIII, such as singular \ird{bró}
and plural \ird{brá}, select cumulative case-number--anchoring endings.
Their long vowels belong to those endings, while the root \ird{br-}
remains constant (§\,\ref{sec:nm-def-allomorph}).

\subsection{Fleeting \ird{e} and reduced stems}\label{sec:vowel-zero}

Alternating nouns such as \ird{Janek}, \ird{perek}, and \ird{obel} have
bound stems \ird{Jank-}, \ird{pérk-}, and \ird{obl-}. Across these alternating classes one
short unstressed \ird{e} in the final stem syllable alternates with zero.
Deletion itself preserves the retained vowel and consonant sequence;
an independently selected quantity grade can supply its long counterpart,
with the high-vowel output of historical \ird{é, ó} in the principal
spoken variety:

\pex
\a \ird{Janek} + \ird{-u} $\rightarrow$ \ird{Janku}
\a \ird{perek} + \ird{-u} $\rightarrow$ \ird{pérku}
\a \ird{obel} + \ird{-u} $\rightarrow$ \ird{oblu}
\a \ird{Merek} + \ird{-í} $\rightarrow$ \ird{Mercí}
\xe

The term \term{fleeting} or \term{weak e} names this alternation, not an
additional phoneme or a reduced vowel obligatorily distinct in pronunciation.
Stable stems retain ordinary \ird{e} in the same morphological environment.
Membership in the alternating class must therefore be acquired with the lexeme;
within that class the relation between full and bound stems is regular. The
analysis locates lexical exceptionality on the morpheme rather than on an
otherwise inaudible vowel distinction, following the representational option
developed for Russian by Gouskova \cite{gouskova2012}. The Iridian class has its own distribution, which is specified by the
paradigms below.

Bound-stem morphology selects the reduced stem before the licensed case-number
endings and, in alternating nouns, before direct anchoring. Selected verbal and
nominalising constructions also admit this relation. The purposive form
\ird{raznit}, for example, selects the reduced stem \ird{razn-} of
\ird{razeň-}; the hard nasal belongs to that lexical bound stem. Vowel loss
alone does not predict the loss of softness.

Deletion is blocked if the remaining consonants cannot be syllabified around
vocalic nuclei without unlicensed consonant loss. Obstruent--sonorant onsets
and sonorant--obstruent coda--onset sequences provide different syllabifications:
\ird{oblu} has \ird{o.blu}, while \ird{pérku} has \ird{pér.ku}. The first
vowel of \ird{perek} remains stressed and cannot delete. Liquids are always
non-syllabic (§\,\ref{sec:vocalic-nuclei}). Ordinary voicing assimilation can
apply to the retained cluster, but consonant deletion is not a productive means
of enabling vowel loss.

New lexical roots and loans have stable vocalism unless assigned to an
alternating class. Native nominal formations with alternating \ird{-ek} can
productively assign the fleeting-vowel pattern to the derived noun; this nominal
suffix is distinct from verbal perfective \ird{-ek}. In longer derivatives,
stem selection applies in the domain of the relevant construction, so an outer
suffix does not automatically delete a vowel inside an already formed stem.

\subsection{Raised grades and vowel quantity}\label{sec:vowel-raising}

Vocalic grades are selected independently of consonantal grades.
G0 is the lexical or plain grade; G1 is strong, selecting the long
correspondent of an eligible short nucleus; G2 is the raised grade, whose
short-vowel correspondences preserve quantity; G3 combines the raised and
strong specifications. The four morphological grades retain their
established assignments and traditional spellings, although their vowel
outputs are only partly distinct in the principal spoken variety: strong \ird{é, ó} have /iː, uː/, as do raised strong
\ird{í, ú}. Conservative varieties retain /eː, oː/ in G1
(Table~\ref{tab:long-vowel-dialects}). The selected nucleus belongs
to the root or stem domain licensed by the construction. It is not every
written vowel in the word. Unless a lexeme specifies another target, the
productive target is the rightmost eligible lexical-root nucleus, excluding
fleeting vowels, prefixes, and suffixes. The \orth{j} of a labial
digraph is not a nucleus. Thus \ird{zemac-} \trsl{run} has two root nuclei and
targets its \ird{a}, giving strong \ird{zemác-}. The forms
\phon{zemácime}{ˈzɛ.maː.tsɪ.mɛ}{progressive run} and
\phon{zemácek}{ˈzɛ.maː.tsɛk}{perfective run} retain initial stress
while lengthening the second root nucleus. Their affricate remains /ts/;
quantity selection supplies no /ts/ $\rightarrow$ /tʃ/ mutation.
The low vowel of \ird{kan-} selects strong \ird{kán-}, with /aː/,
without a height alternation. Thus
\phon{kánek}{ˈkaː.nɛk}{perfective sing} combines root G1 with a short
perfective suffix.
Morphologically selected length and quantity alternations have comparative
analyses in \cite{quantity2023zimmermann} and
\cite{zimmermanntrommer2013}; the four grades and their Iridian assignments
are descriptive specifications here, not consequences of those comparisons.

\begin{table}[ht]
\footnotesize
\caption{Vocalic grades for eligible short mid nuclei: traditional spelling
and principal spoken pronunciation.}\label{tab:vocalic-grades}
\begin{tabularx}{\textwidth}{l L L L L}
\toprule
{\sc input} & {\sc G0 plain} & {\sc G1 strong} & {\sc G2 raised} & {\sc G3 raised strong} \\
\midrule
\ird{e/ě} & \ird{e/ě} /ɛ/ & \ird{é} /iː/ & \ird{i} /ɪ/ & \ird{í} /iː/ \\
\ird{o} & \ird{o} /ɔ/ & \ird{ó} /uː/ & \ird{u} /ʊ/ & \ird{ú} /uː/ \\
\bottomrule
\end{tabularx}
\end{table}

The noun \ird{lobek} selects G3 \ird{lúbk-} throughout its bound paradigm,
including \ird{lúbku}, \ird{lúbcí}, and \ird{lúbce}. By contrast,
\ird{perek} selects G1 \ird{pérk-}, despite sharing
\ird{-u} with \ird{lobek}. Height and strength therefore belong to the
selected stem class, independently of the surface vowel of the ending.

G2 has the following traditional spelling correspondences; the short
series remains a height alternation in all three regional systems, whereas
the long series depends on whether the relevant merger is complete:

\begin{center}
\begin{tabular}{ll}
short: & \ird{e/ě} $\leftrightarrow$ \ird{i}, \quad \ird{o} $\leftrightarrow$ \ird{u} \\
long: & \ird{é} $\leftrightarrow$ \ird{í}, \quad \ird{ó} $\leftrightarrow$ \ird{ú}
\end{tabular}
\end{center}

The short high vowels have their ordinary lax realisations [ɪ, ʊ].
Conservative long /eː, oː/ correspond to raised /iː, uː/, while the
principal variety has /iː, uː/ on both sides of each long spelling
correspondence. In the region with only the back-vowel merger,
\ird{ó}/\ird{ú} coincide at /uː/ but \ird{é}/\ird{í} retain
/eː/ versus /iː/. Since \ird{ě} and \ird{e} represent the same vowel
phoneme, they do not define different height grades. Lexical consonantal
palatality is retained independently of these vowel correspondences: \ird{zděk-} contains /ɟ/, which remains /ɟ/ in
raised \ird{zďíčev}; \ird{pěštel} retains its lexical /pʲ/ in
\ird{pištlí}, where high-front-vowel spelling does not separately mark it.
G2 alone does not lengthen a short root vowel. G3 instead maps short
\ird{e/ě} to \ird{í} and short
\ird{o} to \ird{ú}; lexical long \ird{é, ó} become \ird{í, ú} without
shortening or a further change of vowel identity in the principal variety.
G1 is available for all six short-vowel classes, although their long
outputs no longer form six distinct phonemic counterparts. G2 and G3
require membership in the licensed mid-vowel classes; no further height relation is
inferred for \ird{a, i, u, y}. No grade supplies a nucleus to a cluster.
Strong selection is idempotent: a nucleus already long does not acquire an
additional degree of length. G0 retains a lexical long vowel rather than
turning it into a short vowel. Where a strong construction
has no licensed height alternation, as with the \ird{a} of irrealis
\ird{pjášč-}, strength changes quantity without inventing a new height grade.

For eligible short mid roots, G1 and G3 therefore coincide in their
root-vowel outputs at /iː/ or /uː/, while G0 and G2 retain the contrast
between short mid and short high vowels. In historical long mid-vowel roots,
G0, G1, G2, and G3, where licensed, have the same vowel identity and
quantity after the relevant merger. Grade selection still determines the traditional spelling
and remains independent of consonantal grades and overt endings; coincidence
of the root vowels alone does not establish identity of complete words.
In the principal variety, the synchronic G1 correspondence /ɛ/ to /iː/
or /ɔ/ to /uː/ is stated directly, without an intermediate phoneme
/eː/ or /oː/. The historical and conservative relations explain the
traditional grade terminology. Traditional morphophonological
representations can accordingly retain separate \ird{é, ó} classes,
whose conservative values are /eː, oː/, as reference inputs to a
variety-specific realisation; these classes introduce no additional
phonemic contrasts in the principal variety.

The high front output of \ird{é} conditions ordinary labial
palatalisation and /x/ fronting, so \ird{pérku} has
[ˈpʲiːr.kʊ] in the principal variety. This contextual articulation
applies after morphological grade selection and introduces no further
coronal PAL or velar mutation; a hard coronal before raised \ird{é}
retains its selected consonantal identity.

\pex
\a \ird{lobek} + \ird{-u} $\rightarrow$ \ird{lúbku} \lingcontext{reduced G3 bound stem}
\a \ird{pěštel} + \ird{-í} $\rightarrow$ \ird{pištlí} \lingcontext{reduced stem; short root vowel raises to short i}
\a \ird{zděk-} + \ird{-ěv} $\rightarrow$ \ird{zďíčev} \lingcontext{G3 modal grade; k $\rightarrow$ č; no vowel loss}
\xe

Verbal raising belongs to specified stem--construction combinations, especially
the irrealis. It does not follow from front-vocalic suffixation as such:
\ird{zděcir}, with negative irrealis \ird{-ir}, retains its mid root vowel.
The linker-selected \ird{-ik-} grade of the perfective is suffixal allomorphy,
not an instruction to raise every preceding root vowel
(§\,\ref{sec:pf-dissim}).

Affirmative irrealis assigns its root grade only where the root is directly
exposed. An intervening causative, applicative, or orientation head blocks that
assignment: Mood does not reach back through it. A local exponent can instead
bear strength, as in PFV \ird{-ov-} becoming \ird{-óv-} before affirmative
irrealis. Negative \ird{-ir} selects neither G3 nor a raised root grade.

The nominal raised grades are distinct from the productive assignment
of fleeting-vowel membership in native formations. An unfamiliar noun with a
mid vowel and ending \ird{-u} does not automatically join the raised class.
Quantity grades, including \ird{stráz}/\ird{straz-}, are specified separately
(§\,\ref{sec:comp-lengthening}).

\section{Consonantal alternations}

\subsection{Soft-stem selection}\label{sec:front-softening}

Both front and mixed grades select PAL for their coronal and labial targets:
/t\,d\,n\,l/ gives [c\,ɟ\,ɲ\,ʎ], /s\,z\,ts/ gives /ʃ\,ʒ\,tʃ/,
and labials receive secondary palatalisation. Their different velar
outputs do not alter this shared correspondence set. The standard
spelling writes soft coronals as \ird{tě, dě, ně, lě} before /ɛ/ and
\ird{ť\,ď\,ň\,ľ} elsewhere, with \ird{š\,ž\,č} for sibilant outputs.
Consonantal specifications are retained independently of vowel grades,
including in \ird{pištlí}
(§\,\ref{sec:orthography}). A rhotic
retains /r/ without obligatory secondary palatalisation
(§\,\ref{sec:soft-coronals}).

Consonantal grade interpretation precedes syllabification and coda lateral vocalisation,
so its lateral output [ʎ] remains unaffected by vocalisation of hard /l/.
The secondary palatal articulation of /v/ is realised on prevocalic [ʋʲ]
(§§\,\ref{sec:l-vocalisation} and~\ref{sec:v-realisation}).

The unglossed stems \ird{rat}, \ird{rad}, \ird{ran}, \ird{ral}, and
\ird{rac} below are schematic inputs illustrating the final-consonant
correspondences, without assigned lexical meanings.

\pex
\a \ird{rat} + \ird{-e} $\rightarrow$ \ird{ratě}
\lingcontext{PAL linker selects [c]}
\a \ird{rad} + \ird{-í} $\rightarrow$ \ird{raďí}
\lingcontext{PAL nominal ending selects [ɟ]}
\a \ird{ran} + \ird{-ě} $\rightarrow$ \ird{raně}
\lingcontext{PAL selects /ɲ/; ně spells the selected consonant and /ɛ/}
\a \ird{ral} + \ird{-e} $\rightarrow$ \ird{ralě}
\lingcontext{PAL linker selects [ʎ]}
\a \ird{les} + \ird{-ě} $\rightarrow$ \ird{leše}
\lingcontext{indirect nominal ending; PAL s $\rightarrow$ š}
\a \ird{stráz} + \ird{-ě} $\rightarrow$ \ird{straže}
\lingcontext{short nominal stem grade; PAL z $\rightarrow$ ž}
\a \ird{rac} + \ird{-e} $\rightarrow$ \ird{rače}
\lingcontext{schematic /ts/-final input; PAL selects /tʃ/}
\xe

Already palatal consonants retain their quality. The postalveolars
/ʃ\,ʒ\,tʃ/ retain their articulation, including when selected as outputs
of this productive mapping.
Sibilant replacement in antipassives is governed by its own construction.

Mutation operates on the selected bound stem, so vocalic and consonantal
changes must be distinguished:

\pex
\a \ird{lobek} + \ird{-í} $\rightarrow$ \ird{lúbcí}
\lingcontext{reduced raised stem; front-grade k $\rightarrow$ ts}
\a \ird{Janek} + \ird{-í} $\rightarrow$ \ird{Jancí}
\lingcontext{reduced stem; front-grade k $\rightarrow$ ts}
\a \ird{perek} + \ird{-e} $\rightarrow$ \ird{pérce}
\lingcontext{first vowel retained; second e deletes; k $\rightarrow$ ts}
\xe

The /ts/ in \ird{lúbcí} and \ird{Jancí} remains [ts] before /iː/.
Negative irrealis \ird{-ir} selects
only the restricted velar-front grade; conditional \ird{-ej} selects the
plain grade. Neither is an exception to an automatic front-vowel mutation
rule, since no such rule governs morphological grade selection.

\subsection{Velar fronting and affrication}\label{sec:velar-fronting}

Participating /k/-final stems select /ts/ in the front grade and /tʃ/ in the
mixed grade. These are affricates, written \ird{c} and \ird{č}, rather than
the secondarily palatalised velar [kʲ] of lexical \ird{kě} or the palatal
stop [c] of \ird{tě}. The distinction between secondary
velar palatalisation and affrication is also necessary in descriptions of
Polish \cite{rubach2019}, although the Iridian classes have their own
membership and outputs.

\begin{table}[ht]
	\footnotesize
	\caption{Class-selected velar alternants.}\label{tab:velar-fronting}
	\begin{tabularx}{\textwidth}{L L L}
		\toprule
		{\sc grade} & {\sc output} & {\sc derivation} \\
		\midrule
		Front & /k/ $\rightarrow$ /ts/ & \ird{Marek} + \ird{-í} $\rightarrow$ \ird{Marcí} \\
		& & \ird{ruk} + \ird{-ě} $\rightarrow$ \ird{ruce} \\
		Restricted velar-front & /k/ $\rightarrow$ /ts/ & \ird{zděk-} + \ird{-ir} $\rightarrow$ \ird{zděcir} \\
		Mixed & /k/ $\rightarrow$ /tʃ/ & \ird{zděk-} + \ird{-ěv} $\rightarrow$ \ird{zďíčev} \\
		\bottomrule
	\end{tabularx}
\end{table}

At a single boundary the input /k/ selects its velar alternant once.
The resulting /ts/ does not re-enter the coronal PAL mapping at that same
boundary. Thus \ird{ruk} plus front-grade \ird{-ě} gives /rʊtsɛ/,
spelled \ird{ruce}, while input /ts/ with PAL before /ɛ/ gives /tʃɛ/,
spelled \ird{če}. The spellings \ird{ce} and \ird{če} record the
different consonants selected from /k/ and /ts/. A distinct outer suffix supplies a new mapping only where its construction
independently licenses the selected grade. That mapping targets the exposed
edge of its own input, rather than returning to the original /k/.

The mixed-grade \ird{pjášč-} of \ird{pjašt-} is a listed cluster alternant;
it does not establish a productive /t/ $\rightarrow$ /tʃ/ mapping.
Likewise, the long raised vowel in \ird{zďíčev} belongs to the selected
G3 modal grade, not to the affrication rule. The contrast with \ird{zděcir}
therefore involves independently selected consonantal and vocalic grades.

\noindent\textit{Treatment of h}\quad

The fricative /x/ does not undergo productive categorical mutation to a
sibilant. Contextual fronting gives [ç] before /ɪ, iː/ and before a
palatal consonant in the onset, as in \ird{htě} [çcɛ], while lexical
/xʲ/ has [ç] in \orth{hě}. Neither contextual fronting nor lexical
softness selects a /ʃ/ stem.
Any exceptional lexical alternant must be specified independently.

\subsection{Antipassive stem replacement and augmentation}\label{sec:sibilant-replacement}

The antipassive construction is conventionally labelled \ird{-aš-}, but its
synchronic exponence is class-conditioned stem replacement or augmentation,
together with a vocalic grade where the lexeme specifies one. The label
names this morphological construction; the surface forms are built from
the resulting stem and the overt aspect or other endings that follow it.
A derivation of \ird{pjášček}, for example, selects antipassive
\ird{pjášč-} from \ird{pjašt-}, then combines it with perfective
\ird{-ek}. It does not first concatenate a uniform /aʃ/ sequence and
subsequently delete that sequence without an independent phonological rule.

The correspondences in Table~\ref{tab:sibilant-mutation} act on the
thematic consonant at the selected stem's right edge. They are the
exponents of the antipassive construction, rather than an additional
application of the front or mixed grade. Availability of a consonantal stem
shape alone does not establish that a verb licenses the antipassive
argument frame (§\,\ref{sec:vm-antipassive-desc}).

\begin{table}[ht]
	\footnotesize
	\caption{Antipassive stem replacement and augmentation classes.}\label{tab:sibilant-mutation}
	\begin{tabularx}{0.9\textwidth}{l L L L}
		\toprule
		{\sc class} & {\sc thematic consonants} & {\sc stem operation} & {\sc example} \\
		\midrule
		I    & -t, -k, -c, -č & replace with -č & \ird{pjašt-} $\rightarrow$ \ird{pjášč-} \\
		II   & -s, -š & replace with or retain -š & \ird{tras-} $\rightarrow$ \ird{tráš-} \\
		II-A & -l, -p & append -š & \ird{dal-} $\rightarrow$ \ird{dalš-} \\
		III  & -d, -h, -j, -z, -ž & replace with or retain -ž & \ird{vad-} $\rightarrow$ \ird{váž-} \\
		III-A & all remaining & append -ž & \ird{ščen-} $\rightarrow$ \ird{ščenž-} \\
		\bottomrule
	\end{tabularx}
\end{table}

Classes I, II, and III select replacement stems, with identity where the
input already has the specified final consonant. Classes II-A and III-A
retain the thematic consonant and augment it with a following sibilant.
Replacement classes select the strong root grade independently of their
consonantal operation: \ird{pjášč-}, \ird{tráš-}, and \ird{váž-} have G1.
Additive \ird{dalš-}, \ird{kanž-}, and \ird{lovž-} preserve their
short roots with their listed short aspect endings. Aspect does not reach
through additive Voice to assign G1 to the lexical root.
Surface [g]-bearing loans, where
assigned to this construction, follow the voiced velar/fricative series
of Class~III without establishing a native /g/ phoneme.

Subsequent morphology operates on the exposed edge of the selected
antipassive stem, and the ordinary syllabification, coda, and voicing rules
apply to the resulting form. Shared postalveolar outputs do not collapse
antipassive exponence into PAL: /t/ gives \ird{č} under antipassive
replacement but \ird{ť} under coronal PAL, while /l/ selects
augmentation in \ird{dalš-} rather than the PAL counterpart \ird{ľ}.

\subsection{Cluster repair, epenthesis, and voicing cleanup}\label{sec:cluster-repair}

Derived consonant clusters are retained wherever they can be syllabified around
vocalic nuclei. Weak-\ird{e} is retained when deletion would prevent this
syllabification without unlicensed consonant loss; clusters created by subsequent
morphological operations undergo assimilation or class-licensed simplification, with vowel epenthesis supplying
a nucleus when consonantal repair is insufficient. Liquids remain non-syllabic
throughout these operations (Section~\ref{sec:vocalic-nuclei}).

\begin{table}[ht]
	\footnotesize
	\caption{Cluster retention and repair at morpheme boundaries.}\label{tab:cluster-repair}
	\begin{tabularx}{\textwidth}{l L L}
		\toprule
		{\sc stage} & {\sc strategy} & {\sc effect} \\
		\midrule
		1 & deletion or retention & delete weak \ird{e} only if the remaining vowels support a licit syllabification; otherwise retain it \\
		2 & cluster retention & preserve derived consonant sequences that are already licit \\
		3 & assimilation or simplification & apply ordinary assimilation or a lexically licensed simplification; do not delete consonants to enable vowel loss \\
		4 & vowel epenthesis & insert a short vowel if a later cluster cannot be repaired by consonantal means \\
		\bottomrule
	\end{tabularx}
\end{table}

Four repairs are central in the modern language. First, obstruent clusters obey
regressive voicing assimilation, including across morpheme boundaries:

\pex
\a \ird{obez + -ba} $\rightarrow$ \ird{obezba}, phonetically [ˈʔɔ.bɛz.ba]
\a \ird{ot + -ba} $\rightarrow$ [ˈʔɔd.ba]
\a \ird{vrad + -ka} $\rightarrow$ [vratka]
\xe

Second, the obstruents remaining after coda vocalisation undergo final
devoicing, while [w] from /l/ or /v/ remains voiced
(§\,\ref{sec:lenition}). Third, weak-vowel stems retain enough
vocalism for their liquids to occupy syllable margins. In the following
constructed derivations, only the second stem vowel deletes:

\pex
\a \ird{berel} + \ird{-a} $\rightarrow$ \ird{berla}
\lingcontext{constructed plain-grade suffixation}
\a \ird{perek} + \ird{-u} $\rightarrow$ \ird{pérku}
\xe

Both outputs contain two syllables, each with a vowel nucleus. The sequences
\ird{rl} and \ird{rk} span coda--onset boundaries, with \ird{r} in the coda
of the first syllable.

The same syllabification precedes coda /l/ vocalisation. An onset lateral
in \ird{berla} or \ird{oblu} remains [l], whereas a hard lateral assigned
to a coda surfaces as [w]. Vocalisation does not license additional fleeting-vowel
deletion, consonant loss, or syllabic consonants
(§\,\ref{sec:l-vocalisation}).

Fourth, a cluster introduced by later consonantal morphology may require
assimilation or class-licensed simplification before vowel epenthesis. Such
repairs do not license an otherwise blocked fleeting-vowel deletion.
Anticipatory palatal coarticulation may affect the consonant immediately
preceding a soft consonant within an onset, especially a coronal. This
phonetic tendency supplies no categorical PAL mapping or additional stem
allomorph. Productive mutation targets the consonant at the stem's right
edge; listed cluster alternants such as \ird{pjašt-}/\ird{pjášč-}
belong to their construction and do not establish unrestricted leftward
spreading. If no consonantal repair succeeds, a short
\ird{e} /ɛ/ is inserted and written as an ordinary vowel. This vowel supplies
a syllable nucleus and is subject to the regular stress rule. It introduces no
morphological softening trigger, since soft-stem selection is conditioned by
suffix class rather than by an inserted vowel.

Vowel retention preserves a segment already present in the stem, whereas
epenthesis adds a vowel to repair a derived cluster. The latter is analysed
here as a phonological segment with nuclear status, following Hall's distinction
between epenthetic vowels and non-nuclear phonetic vowel intrusion
\cite{hall2006}. Neither process licenses a syllabic
liquid, and neither entails a new vowel phoneme beyond the inventory in
Table~\ref{tab:vowels}.

\subsection{Quantity grades}\label{sec:comp-lengthening}

Quantity is a morphological contrast independent of stress and consonantal
mutation. Native short Class~I nouns select G0 in the singular and G1 throughout
the plural, with short transitive plural \ird{-ov}; loans and invariant-long
stems have explicitly specified subclasses. Class~II distinguishes G1
\ird{pérk-} from G3 \ird{lúbk-}, while retaining the case endings of each
subclass: singular indirect \ird{pérce} and \ird{lúbcí} are distinct cells.
Class~III \ird{trava} selects G0 in the singular and G1 throughout the plural,
including \ird{tráv} and plural \ird{-am}. The restricted alternating
\ird{stráz}/\ird{straz-} subtype retains its own distribution.
Outer anchoring preserves the selected inner quantity and protected long
endings: \ird{pércé} and \ird{lúbcé} are not shortened. Lexical long vowels,
as in \ird{rán-}, are not globally shortened by G0.

\pex
\a \ird{vad-} $\xrightarrow{\mathrm{ANTIP}}$ \ird{váž-} \lingcontext{replacement stem with G1}
\a \ird{dal-} $\xrightarrow{\mathrm{ANTIP}}$ \ird{dalš-} \lingcontext{stem augmentation; quantity preserved}
\xe

The short G2 correspondence preserves quantity, whereas strong selection
assigns a long output to an eligible short nucleus. The complete mergers
of historical long mid vowels preserve length and apply independently of
these morphological selections; neither grade selection nor merger is a
general response to a dense consonant cluster.

\noindent\textit{Verbal strength and locality.}\quad
Class~A short roots place strength in the root in progressive and perfective
forms: \ird{zemaca}, \ird{zemácime}, \ird{zemacáž}, \ird{zemácek},
\ird{zemacáv}. Class~B \ird{pjašt-} instead places it in the aspect ending:
\ird{pjašta}, \ird{pjaštíme}, \ird{pjaštáž}, \ird{pjaštník},
\ird{pjaštnáv}. Class~C lexical long stems retain their root quantity with
short inner endings, including contemplative \ird{-až}. These classes do
not follow from the position of stress.

New strength introduced by Voice, causative, or applicative belongs to the
local derived stem. Compact causative gives \ird{pjáštr-} and \ird{vádr-};
expanded causative gives \ird{ščenér-} where the preceding stem lacks
quantity, but short \ird{-er-} follows already strong \ird{pjášč-}.
Mutating \ird{vaďér-} combines PAL with suffixal strength. Protected
applicative \ird{-éb-} retains its long vowel in \ird{pjaštébr-} and
blocks outer causative root lengthening. Immediately following aspect endings
are short: \ird{-ime, -ek/-ik, -av, -až}. This restricted rhythmic
allomorphy concerns the adjacent inner Asp exponent; it is not a
one-long-vowel-per-word rule and does not shorten protected anchoring or
independent lexical quantities, directive Mood, or outer nominalisers
\ird{-ou} and \ird{-ový}. Slovak rhythmic shortening supplies a
comparative case of quantity interaction with language-specific exceptions
\cite{pukanec2016}; it does not determine these Iridian allomorphs or their
morphological domains. An outer weak selector does not reset the
quantity of an intervening strong derived stem.

\section{Orthographic realisation}

\subsection{Orthography and the soft consonant letters}\label{sec:orthography}

Standard Iridian represents consonantal outputs after lexical and
morphological selection. Before /ɛ/, softness in the coronal, labial, and
restricted dorsal classes is marked by \orth{ě}; soft coronals otherwise
use \ird{ť\,ď\,ň\,ľ}. Before pronounced \ird{i, í}, labial softness
requires no additional marker, although a lexical specification remains
part of the stem. In the principal variety, \ird{é} also conditions
contextual labial palatalisation, while retaining its traditional spelling,
as in \ird{pérku} [ˈpʲiːr.kʊ]. The long spellings \ird{é, ó}
continue to distinguish their historical and morphological classes from
\ird{í, ú}, although the paired vowels coincide in pronunciation. A vowel
letter supplies no additional PAL operation.
The principal consonant conventions are:

\begin{itemize}
	\item \ird{č}, \ird{š}, \ird{ž} spell postalveolar phonemes, whether
	lexical or derived through coronal PAL, a mixed-grade velar output, or antipassive replacement.
	\item Bare \ird{c} spells /ts/, including before high front vowels;
	\ird{č} spells /tʃ/. Neither letter denotes the palatal stop [c].
	\item The phoneme /x/ is written \ird{h} in every position, including
	word-finally. The same letter represents a single velar fricative throughout
	the paradigm; it does not denote a consonant cluster.
\end{itemize}

\subsection{The letter ě and orthographic normalisation}\label{sec:e-hook}

Orthographic normalisation records the selected consonantal output together
with the following vowel. Lexical and morphologically selected softness
receive the same spelling before /ɛ/, while plain consonants retain
\ird{e}. The sequences in Table~\ref{tab:e-norm} illustrate spelling
correspondences rather than additional lexical entries.

\begin{table}[ht]
	\footnotesize
	\caption{Spellings before /ɛ/, by consonantal output.}\label{tab:e-norm}\label{tab:ortho-e-hook}
	\begin{tabularx}{\textwidth}{L L L}
		\toprule
		{\sc consonantal output} & {\sc plain + /ɛ/} & {\sc soft + /ɛ/} \\
		\midrule
		\ipagloss{[t] : [c]} & te & tě \\
		\ipagloss{[d] : [ɟ]} & de & dě \\
		\ipagloss{[n] : [ɲ]} & ne & ně \\
		\ipagloss{[l] : [ʎ]} & le & lě \\
		\ipagloss{[p] : [pʲ]} & pe & pě \\
		\ipagloss{[b] : [bʲ]} & be & bě \\
		\ipagloss{[m] : [mʲ]} & me & mě \\
		\ipagloss{[ʋ] : [ʋʲ]} & ve & vě \\
		\ipagloss{[k] : [kʲ]} & ke & kě \\
		\ipagloss{[x] : [ç]} & he & hě \\
		\bottomrule
	\end{tabularx}
\end{table}

Thus linker \ird{-e} gives \ird{ratě} from schematic \ird{rat},
while indirect \ird{-ě} gives \ird{raně} from schematic \ird{ran}.
Lexical \ird{tězeň}/\ird{tězn-} follows the same convention, with
linked \ird{tězně}; \ird{zděk-}, \ird{blětébov-}, and the lexical
forms \ird{nadě, hladě, nět, mně, jeně, Štěpánov, kaděc, hloutěc, priděc}
likewise contain soft coronals wherever \ird{ě} follows a coronal letter.
Their spelling assigns no independent derivational boundary or productive
morphological analysis. Lexical \ird{kěda} [ˈkʲɛ.da] and \ird{těda}
[ˈcɛ.da] retain different places of closure despite sharing \ird{ě}.

The postalveolars and plain sibilants retain their own letters before
\ird{e}: \ird{še, že, če} contrast with \ird{se, ze, ce}.
In particular, \ird{ce} [tsɛ] and \ird{če} [tʃɛ] preserve the
affricate contrast.
Redundant \ird{*ťě, *ďě, *ňě, *ľě} and \ird{*čě, *šě, *žě}
are excluded.
After unchanged /r/ and independent /j/, /ɛ/ is written \ird{e}, as in
schematic \ird{re, je}. Accordingly, causative-final \ird{r} takes
front-grade \ird{-ěc} and mixed-grade \ird{-ěv} with surface \ird{e}:
\ird{pjáštrec}, \ird{pjáštrev}. Neither suffix mutates the lexical
root across the intervening causative.

\subsection{The letter ř}\label{sec:r-hook}

The letter \ird{ř} is a lexical spelling of /ʒ/ in learned and
etymological vocabulary. It introduces no contrast with \ird{ž}, has the
same devoiced realisation [ʃ] word-finally, and does not represent either
a palatalised rhotic or Czech /r̝/. It is preserved where specified by a
lexical entry, as in \ird{knář-}, but no productive native inflection
changes /r/ to \ird{ř}. A phonetic [rʲ] retains rhotic articulation and
does not receive this spelling.

\subsection{Lexical class membership}

A lexeme's stored stem relations and its declension class jointly determine the
form selected by each construction. Expanded endings such as \ird{-ení}
and animate-like \ird{-ovi} do not collapse this distinction: their
consonantal grades and the stem's vocalic grade remain separate
specifications. Listed irregular forms supplement the productive mappings;
they do not supply a general phonological rule for unfamiliar words.

\subsection{Perfective suffix grades}\label{sec:pf-dissim}

Free perfective \ird{-ek} corresponds to bound \ird{-ik-} before
the linker \ird{-e} and participant nominaliser \ird{-ou}:
\ird{pjášček}, \ird{pjáščice}, \ird{pjáščikou}. The linker also triggers
\ird{k $\rightarrow$ c}, while \ird{-ou} preserves the velar. These
constructions select the suffixal grade without raising the lexical root.
Class~A assigns root strength, giving \ird{ščen-} $\rightarrow$
\ird{ščének}; a front mid root vowel does not itself select \ird{-ik}.

Class~B has free intransitive \ird{-ék} and bound \ird{-ík-}, and
transitive \ird{-ník} retains its quantity when bound. The long
\ird{-ék}/\ird{-ík-} alternation is retained in spelling, but both
have /iː/ in the principal variety; conservative pronunciation retains
/eː/ versus /iː/. Their consonantal selections remain independently
specified. Strong local
Voice/Caus/Appl and Class~C instead take the short adjacent endings.
Perfective \ird{-ik} selected by stative \ird{-n-} remains a
construction-specific form; it is not the default of every front-mid root.

\subsection{Rule interactions}\label{sec:morphophon-interactions}

Stem-grade selection and suffixal mutation combine in the following forms:

\pex
\a \ird{lobek} + \ird{-u} $\rightarrow$ \ird{lúbku}
\lingcontext{suffix class selection; weak-\ird{e} deletion; back-vowel raising; no front-vocalic mutation}
\a \ird{Janek} + \ird{-í} $\rightarrow$ \ird{Jancí}
\lingcontext{weak-\ird{e} deletion; front-grade k $\rightarrow$ ts; [ts] retained before í}
\a \ird{perek} + \ird{-e} $\rightarrow$ \ird{pérce}
\lingcontext{CeLeC stem retains its first vowel and loses the second; velar fronting k $\rightarrow$ c follows}
\a \ird{ruk} + \ird{-ě} $\rightarrow$ \ird{ruce}
\lingcontext{front-grade k $\rightarrow$ ts yields /tsɛ/, written ce; no second PAL application}
\a \ird{louž} + \ird{-ení} $\rightarrow$ \ird{loužení}
\lingcontext{expanded front-vocalic indirect ending; stem-final ž retains its postalveolar articulation}
\a \ird{zděk-} + \ird{-ěv} $\rightarrow$ \ird{zďíčev}
\lingcontext{G3 raised strong modal grade; k $\rightarrow$ č; no vowel deletion}
\a \ird{zděk-} + \ird{-ir} $\rightarrow$ \ird{zděcir}
\lingcontext{restricted velar-front grade; k $\rightarrow$ ts; no coronal PAL or root raising}
\a \ird{pjašt-} $\xrightarrow{\mathrm{ANTIP}}$ \ird{pjášč-};
\ird{pjášč-} + \ird{-ek} $\rightarrow$ \ird{pjášček}
\lingcontext{Class I replacement expones antipassive; overt perfective suffixation follows}
\a \ird{pjášček} + \ird{-e} $\rightarrow$ \ird{pjáščice}
\lingcontext{linker selects perfective -ik-; final k $\rightarrow$ c; the earlier č remains unchanged}
\xe

These derivations are cited throughout Chapters~\ref{ch:verbs}, \ref{ch:nouns},
and~\ref{ch:complex}.

\section{Operational summary}

A construction first selects the full or bound stem and any independently
specified height or quantity grade. Fleeting-vowel reduction preserves the
stressed vowel and enough vocalism to syllabify all liquids at syllable
margins. Each selecting exponent licenses one original-input mapping at
its own boundary (§\,\ref{sec:boundary-mapping}). Plain grade preserves
the input consonant; front selects coronal/labial PAL and velar /ts/;
mixed selects the same PAL correspondences and velar /tʃ/; restricted
velar-front selects /ts/ from /k/ without PAL. These four licensed
combinations exhaust the consonantal grade system. Independently, vocalic
G0--G3 select lexical, strong, raised, and raised strong nuclei in licensed
local domains. G0 preserves lexical quantity; it is not a general shortening
operation. Restricted short inner endings do not affect protected outer
derivation or anchoring.

The /ts/ output of front-grade /k/ does not become /tʃ/ at that same
boundary. Another mapping requires a distinct, independently licensed
outer construction and targets its own exposed input edge. Antipassive
replacement or augmentation and listed vocalic or cluster allomorphy
are selected separately. The conventional label \ird{-aš-} does not
insert an invariant segmental suffix in a replacement-stem derivation.

After lexical syllabification and allomorph selection, restricted liaison
of negative \ird{z-} supplies an onset to its vowel-initial host;
word-edge glottal insertion applies where the initial onset remains empty
(§§\,\ref{sec:glottal-onset} and~\ref{sec:neg-allomorphy}). Neither
operation adds a nucleus or changes lexical quantity.
Syllabification and licensed repair establish the syllables on which initial
lexical stress is assigned (§\,\ref{sec:stress}); vowel quantity remains
independent. Principal spoken \ird{é, ó} have /iː, uː/ throughout
lexical and grammatical material, preserving quantity while merging the
long mid-vowel outputs with those of \ird{í, ú}. Contextual labial
palatalisation and /x/ fronting apply to the resulting high front vowel
without selecting a further morphological consonantal grade.
Pronounced /ɪ\,iː/ does not neutralise the two sibilant
series. The letters \orth{i, í} always supply vowel nuclei, whereas the
\orth{j} of a restricted labial digraph supplies no independent segment.
Orthographic normalisation writes soft coronals as \ird{tě, dě, ně, lě}
before /ɛ/ and \ird{ť\,ď\,ň\,ľ} elsewhere, retains \ird{š\,ž\,č},
distinguishes \ird{ce/če}, and preserves lexical \ird{kě} for /kʲɛ/,
distinct from /cɛ/ \ird{tě}, as well as lexical \ird{hě} and learned \ird{ř}.
Neither front-vowel spelling nor an
orthographic normalisation rule supplies an additional morphological trigger.
